% ---------------------------------------------------------------------------
% Corrigenda note -- friends of 10 and 20.
% Compile: pdflatex errata_friends_of_10.tex  (twice)
% All numerical claims are reproduced by experiments/verifica_erratas.py
% ---------------------------------------------------------------------------
\documentclass[11pt,a4paper]{article}

\usepackage[T1]{fontenc}
\usepackage[utf8]{inputenc}
\usepackage{amsmath,amssymb,amsthm}
\usepackage[margin=1in]{geometry}
\usepackage{hyperref}
\hypersetup{colorlinks=true,linkcolor=blue,citecolor=blue,urlcolor=blue}

\theoremstyle{plain}
\newtheorem{theorem}{Theorem}[section]
\newtheorem{proposition}[theorem]{Proposition}
\newtheorem{lemma}[theorem]{Lemma}
\newtheorem{corollary}[theorem]{Corollary}
\theoremstyle{definition}
\newtheorem{remark}[theorem]{Remark}

\newcommand{\NN}{\mathbb{N}}
\newcommand{\ZZ}{\mathbb{Z}}
\newcommand{\QQ}{\mathbb{Q}}
\newcommand{\Iab}{I}
\DeclareMathOperator{\lcm}{lcm}

\title{Corrigenda and remarks on some recent results\\ concerning friends of $10$ and $20$}

\author{Caio Comitre Rossi\thanks{Independent researcher. E-mail:
\texttt{kio199@gmail.com}. Verification code:
\url{https://github.com/caioross/solitary10} (see the note in Section~\ref{sec:verif}).}}

\date{\today}

\begin{document}
\maketitle

\begin{abstract}
We record four corrections to recent work on the open problem of whether $10$ is a
solitary number. (i) A combinatorial lemma on partitions used in
\cite{CMM10} (Lemma 2.3) and in \cite{CMM20} (Lemma 23) is false as stated: the
inequality is strict there, while equality holds for the all-ones partition; we give
the corrected statement together with its equality case, and observe that the lemmas
deduced from it (\cite[Lemma 3.6]{CMM10}, \cite[Lemma 24]{CMM20}) are true as stated
but require the corrected version, since their minima are \emph{attained} exactly in
the equality case. (ii) The two inequalities displayed in \cite[Remark 3.7]{CMM10}
are off by one: what follows from \cite[Theorem 1.10]{CMM10} is
$\Omega(m)\ge\omega(N)+2a-2$ and $\Omega(m)\ge\omega(m)+2a-1$; consequently
\cite[Corollary 1.11]{CMM10} should read $N<5\cdot 6^{(2^{K-2a+2}-1)^{2}}$. We also
show that the counting argument proving \cite[Theorem 1.10]{CMM10} is optimal, so the
missing unit cannot be recovered from that argument. (iii) The main theorem of
\cite{M10} yields upper bounds for the $r$-th smallest prime divisor of a friend of
$10$ only for $2\le r\le 5$, and not for every prime divisor as the title and
abstract state: the admissibility threshold in its hypothesis becomes negative
exactly at $r=6$. We make the effective range explicit and record the bound for
$q_{5}$, which is available in \cite{M10} but not stated there. (iv) A misprinted
abundancy factor in Case 12 of the proof of \cite[Theorem 1.2]{CMM10}. None of the
main theorems of \cite{CMM10,CMM20,M10} is affected. All computations reported here
are exact rational computations and are reproducible from the code cited in
Section~\ref{sec:verif}.
\end{abstract}

\section{Notation and scope}

For $n\in\ZZ^{+}$ let $\sigma(n)$ denote the sum of the positive divisors of $n$ and
$\Iab(n)=\sigma(n)/n$ the abundancy index of $n$; $\omega(n)$ and $\Omega(n)$ denote
the number of distinct prime divisors of $n$ and the number of prime divisors counted
with multiplicity, respectively, and $v_{q}(n)$ is the $q$-adic valuation of $n$. A
positive integer $n\ne 10$ is a \emph{friend of $10$} if $\Iab(n)=\Iab(10)=9/5$;
friends of $20$ are defined analogously with $\Iab(20)=21/10$. Whether $10$ (or $20$)
has a friend is open; see \cite{Ward,Thackeray,MM,CMM10,CMM20,M10}.

Throughout, $p_{i}$ denotes the $i$-th prime ($p_{1}=2$), and $q_{r}$ denotes the
$r$-th smallest prime divisor of the hypothetical friend under consideration.

Two remarks on the logical status of what follows. First, since no friend of $10$ is
known, a statement of the form ``if $N$ is a friend of $10$ then $P(N)$'' cannot be
refuted by exhibiting a counterexample; accordingly, except in
Section~\ref{sec:partition} --- where the statement in question is a self-contained
assertion about partitions and is genuinely false --- we do \emph{not} claim that the
statements we discuss are false. We claim that the arguments given do not establish
them, and we state what the arguments do establish. Second, the four items below are
local; the main results of \cite{CMM10,CMM20,M10}, in particular $\omega(N)\ge 7$
for a friend $N$ of $10$ \cite[Theorem 1.2]{CMM10}, are untouched, as are the bounds
of \cite{MM} and the bound $\omega(N)\ge 10$ of \cite{Thackeray}.

\section{A partition lemma}\label{sec:partition}

The following statement appears as \cite[Lemma 2.3]{CMM10} and, verbatim, as
\cite[Lemma 23]{CMM20}.

\begin{lemma}[as stated in {\cite[Lemma 2.3]{CMM10}}, {\cite[Lemma 23]{CMM20}}]
\label{lem:orig}
Let $(c_{1},\dots,c_{k})$ be any partition of $n$, i.e.\ $\sum_{i=1}^{k}c_{i}=n$ with
$1\le k\le n$ and all $c_{i}\in\NN$. Then for any integer $a>e$,
\[
  a\,n<\sum_{i=1}^{k}a^{c_{i}}.
\]
\end{lemma}

This is false. The proof deduces $ac_{i}<a^{c_{i}}$ for every $c_{i}\ge 1$ from the
fact that $\psi(x)=ax-a^{x}$ is strictly decreasing on $[1,\infty)$ for $a>e$
(\cite[Lemma 2.2]{CMM10}, \cite[Lemma 22]{CMM20}); but $\psi(1)=0$, so strict
decrease gives $ac<a^{c}$ only for $c>1$, with equality at $c=1$. Taking the all-ones
partition $c_{1}=\dots=c_{n}=1$ (so $k=n$) yields
\[
  \sum_{i=1}^{n}a^{c_{i}}=na=a\,n ,
\]
so the asserted strict inequality fails for every $n\ge 1$ and every integer $a>e$;
the smallest instance is $n=1$, $a=3$. The correct statement is the following.

\begin{lemma}\label{lem:fix}
Let $(c_{1},\dots,c_{k})$ be a partition of $n$ and let $a>e$ be an integer. Then
\[
  a\,n\le\sum_{i=1}^{k}a^{c_{i}},
\]
with equality if and only if $c_{i}=1$ for all $i$.
\end{lemma}

\begin{proof}
With $\psi(x)=ax-a^{x}$ we have $\psi'(x)=a-a^{x}\log a<0$ on $[1,\infty)$ for $a>e$,
hence $\psi(x)\le\psi(1)=0$ on $[1,\infty)$ with equality only at $x=1$; that is,
$ac\le a^{c}$ for every integer $c\ge 1$, with equality if and only if $c=1$. Summing
over $i$ gives $an=a\sum_{i}c_{i}\le\sum_{i}a^{c_{i}}$, and equality holds if and
only if it holds termwise, i.e.\ if and only if every $c_{i}$ equals $1$.
\end{proof}

The lemma is used, in both papers, only through the following consequence. Recall
from \cite{CMM10,CMM20} the sets
\[
  \mathcal{A}_{n,a}(r):=\Bigl\{\sum_{i=1}^{r}a^{c_{i}}-r\ :\
    \sum_{i=1}^{r}c_{i}=n,\ c_{i}\in\NN\Bigr\},
  \qquad
  \mathcal{L}_{n,a}:=\bigcup_{r=1}^{n}\mathcal{A}_{n,a}(r).
\]

\begin{lemma}[{\cite[Lemma 3.6]{CMM10}}, {\cite[Lemma 24]{CMM20}}]\label{lem:min}
The minimum element of $\mathcal{L}_{2a-1,5}$ is $8a-4$.
\end{lemma}

The \emph{statement} is correct, but the proof given is not, and the defect is
visible inside the papers themselves: the proof concludes with the strict inequality
$4(2a-1)<\sum_{i}5^{c_{i}}-k$, which asserts that $8a-4$ is a strict lower bound and
therefore \emph{not} a minimum, while the note immediately following the proof states
correctly that the minimum lies in $\mathcal{A}_{2a-1,5}(2a-1)$, i.e.\ that it is
attained at the all-ones partition. Lemma~\ref{lem:fix} repairs this at once: for a
partition of $n=2a-1$ into $k\le 2a-1$ parts,
\[
  \sum_{i=1}^{k}5^{c_{i}}-k\ \ge\ 5(2a-1)-(2a-1)\ =\ 4(2a-1)\ =\ 8a-4 ,
\]
with equality exactly for $k=2a-1$ and all $c_{i}=1$. Hence the minimum is $8a-4$ and
it is attained. Nothing further in \cite{CMM10,CMM20} changes.

\section{An off-by-one in \cite[Remark 3.7]{CMM10} and its effect on
\cite[Corollary 1.11]{CMM10}}\label{sec:offbyone}

Let $N=5^{2a}m^{2}$ be a friend of $10$, so that $5^{2a}\,\|\,N$ and $5\nmid m$; in
particular $\omega(N)=\omega(m)+1$. \cite[Theorem 1.10]{CMM10} states
\begin{equation}\label{eq:110}
  \Omega(N)\ \ge\ 2\omega(N)+6a-4 ,
\end{equation}
and its proof is correct once Lemma~\ref{lem:min} is proved as above.
\cite[Remark 3.7]{CMM10} then displays
\[
  \Omega(5^{2a})+\Omega(m^{2})=2a+2\Omega(m)\ \ge\ 2\omega(N)+6a-4
\]
and concludes
\begin{equation}\label{eq:7}
  \Omega(m)\ \ge\ \omega(N)+2a-1 \tag{7}
\end{equation}
and, from \eqref{eq:7},
\begin{equation}\label{eq:8}
  \Omega(m)\ \ge\ \omega(m)+2a. \tag{8}
\end{equation}
The displayed inequality gives $2\Omega(m)\ge 2\omega(N)+4a-4$, i.e.
\begin{equation}\label{eq:7fix}
  \Omega(m)\ \ge\ \omega(N)+2a-2 ,\tag{7$'$}
\end{equation}
which is one unit weaker than \eqref{eq:7}. No parity gain is available: both sides
of $2a+2\Omega(m)\ge 2\omega(N)+6a-4$ are even, so \eqref{eq:7fix} is exactly the
integer consequence of \eqref{eq:110}. Substituting $\omega(N)=\omega(m)+1$ in
\eqref{eq:7fix} gives
\begin{equation}\label{eq:8fix}
  \Omega(m)\ \ge\ \omega(m)+2a-1 ,\tag{8$'$}
\end{equation}
again one unit weaker than \eqref{eq:8}. Since \cite[Corollary 1.11]{CMM10} is
obtained by feeding $\Omega(m)\le K$ into \eqref{eq:7}, its proof yields
$\omega(N)\le K-2a+2$ rather than $\omega(N)\le K-2a+1$, and therefore, by
Nielsen's bound \cite[Lemma 2.5]{CMM10} for odd $9/5$-perfect numbers:

\begin{corollary}[corrected form of {\cite[Corollary 1.11]{CMM10}}]\label{cor:111}
If $N=5^{2a}m^{2}$ is a friend of $10$ with $a,m\in\NN$ and $\Omega(m)\le K$, then
\[
  N<5\cdot 6^{\left(2^{\,\omega(N)}-1\right)^{2}}
   <5\cdot 6^{\left(2^{\,K-2a+2}-1\right)^{2}} .
\]
\end{corollary}

The corrected statement of \cite[Corollary 1.11]{CMM10} is quoted, in its original
form, in \cite[\S1]{M20div}; the same replacement applies there.

It is natural to ask whether \eqref{eq:7} can be recovered by strengthening
\eqref{eq:110} to $\Omega(N)\ge 2\omega(N)+6a-2$ (the least improvement compatible
with the parity of $\Omega(N)=2a+2\Omega(m)$). The counting used in \cite{CMM10} does
not allow it, and this is not a matter of a lossy estimate: it is exactly optimal.

\begin{proposition}\label{prop:tight}
Fix $a\ge 1$ and $s\ge 2a$. Consider the constraints used in the proof of
\cite[Theorem 1.10]{CMM10} for $N=5^{2a}\prod_{i=1}^{s-1}p_{i}^{2\gamma_{i}}$ with
$\omega(N)=s$, namely: there is a set $T\subseteq\{1,\dots,s-1\}$, $|T|=t\ge 1$, and
integers $v_{i}\ge 1$ $(i\in T)$ with $\sum_{i\in T}v_{i}=2a-1$, such that
$2\gamma_{i}+1\equiv 0\pmod{5^{v_{i}}}$ for $i\in T$ --- whence
$2\gamma_{i}\ge 5^{v_{i}}-1$ --- and $2\gamma_{i}\ge 2$ for $i\notin T$. Then
\[
  \min\ \Omega(N)\;=\;2a+\sum_{i\in T}2\gamma_{i}+\sum_{i\notin T}2\gamma_{i}
  \;=\;2s+6a-4 ,
\]
the minimum being attained by $t=2a-1$, $v_{i}=1$ and $2\gamma_{i}=4$ for $i\in T$,
and $2\gamma_{i}=2$ for the remaining $s-2a$ primes.
\end{proposition}

\begin{proof}
Under the stated constraints,
\[
  \Omega(N)\ \ge\ 2a+\sum_{i\in T}\bigl(5^{v_{i}}-1\bigr)+2(s-1-t)
  \ =\ 2(s-1)+2a+\sum_{i\in T}\bigl(5^{v_{i}}-3\bigr).
\]
For every integer $v\ge 1$ one has $5^{v}-3\ge 2v$, with equality if and only if
$v=1$ (indeed $5^{1}-3=2$, and $5^{v}\ge 2v+3$ is immediate for $v\ge 2$). Hence
\[
  \sum_{i\in T}\bigl(5^{v_{i}}-3\bigr)\ \ge\ 2\sum_{i\in T}v_{i}\ =\ 2(2a-1)\ =\ 4a-2 ,
\]
with equality if and only if $t=2a-1$ and $v_{i}=1$ for all $i\in T$. This gives
$\Omega(N)\ge 2(s-1)+2a+4a-2=2s+6a-4$. The displayed configuration satisfies every
constraint (each $i\in T$ has $2\gamma_{i}=4$, so $2\gamma_{i}+1=5\equiv 0
\pmod{5}$) and attains the bound, and it requires $t=2a-1\le s-1$, i.e.\ $s\ge 2a$.
\end{proof}

Consequently \eqref{eq:110} is the exact optimum of the relaxation used to prove it,
and \eqref{eq:7} does not follow from \cite[Theorem 1.10]{CMM10}; whether \eqref{eq:7}
is true is left open here. We note that the analogous passage in the companion
paper \cite{CMM20} is carried out correctly: there
$\Omega(N)\ge 2\omega(N)+6a-5$ with $N=2\cdot 5^{2a}m^{2}$ gives
$1+2a+2\Omega(m)\ge 2\omega(N)+6a-5$, hence $\Omega(m)\ge\omega(N)+2a-3$, which is
the exact integer consequence.

\section{The effective range of \cite[Theorem 1.2]{M10}}\label{sec:range}

\cite[Theorem 1.2]{M10} states: if $n$ is a friend of $10$ and $q_{r}$ ($r\ge 2$) is
its $r$-th smallest prime divisor, then $q_{r}<L\{\log L+2\log(\log L)\}$, where
$L=\lceil\mathcal{A}\omega(n)/\mathcal{B}\rceil$ and $\mathcal{A},\mathcal{B}$ are
coprime positive integers with $\mathcal{A}/\mathcal{B}\in\QQ^{+}\setminus\ZZ^{+}$,
$\mathcal{A}\mathcal{B}(r-2)+2\mathcal{A}+\mathcal{B}>\mathcal{B}^{2}$ and
\begin{equation}\label{eq:cond1}
  \frac{\mathcal{A}}{\mathcal{B}}\;>\;
  \frac{1}{\;\frac{36}{25}\prod_{4\le i\le r+1}\bigl(1-\frac1{p_{i}}\bigr)-1\;} .
\end{equation}
Write
\[
  X_{r}:=\frac{36}{25}\prod_{4\le i\le r+1}\Bigl(1-\frac1{p_{i}}\Bigr),
  \qquad
  F_{r}:=\frac54\prod_{4\le j\le r+1}\frac{p_{j}}{p_{j}-1},
\]
the empty products being $1$; thus $F_{r}X_{r}=9/5$ for every $r\ge 2$, and
\eqref{eq:cond1} reads $\mathcal{A}/\mathcal{B}>1/(X_{r}-1)$. (That this is the
intended reading is confirmed by the instantiations in \cite{M10}: for $r=2,3,4$ the
thresholds used there are $25/11$, $175/41$ and $385/47$, which are exactly
$1/(X_{r}-1)$.)

The final step of the proof in \cite{M10} is
\[
  \Iab(n)<F_{r}\Bigl(1+\frac{\mathcal{B}}{\mathcal{A}}\Bigr)
  <F_{r}X_{r}=\frac95 ,
\]
so it requires $1+\mathcal{B}/\mathcal{A}<X_{r}$. Since
$1+\mathcal{B}/\mathcal{A}>1$ for positive $\mathcal{A},\mathcal{B}$, this is possible
if and only if $X_{r}>1$. The sequence $(X_{r})_{r\ge 2}$ is strictly decreasing, and

\begin{proposition}\label{prop:Xr}
$X_{r}>1$ if and only if $2\le r\le 5$. Explicitly,
\[
  X_{2}=\tfrac{36}{25},\quad X_{3}=\tfrac{216}{175},\quad
  X_{4}=\tfrac{432}{385},\quad X_{5}=\tfrac{5184}{5005},\quad
  X_{6}=\tfrac{82944}{85085}<1 ,
\]
and $X_{r}<X_{6}<1$ for $r>6$. Equivalently, $F_{r}\ge 9/5$ for $r\ge 6$.
\end{proposition}

\begin{proof}
The displayed values are exact rational evaluations, and $X_{r+1}=X_{r}(1-1/p_{r+2})
<X_{r}$. The last sentence follows from $F_{r}X_{r}=9/5$.
\end{proof}

Hence, for $r\ge 6$: under the intended reading of \eqref{eq:cond1} the right-hand
side is negative and no admissible pair $(\mathcal{A},\mathcal{B})$ of positive
integers satisfies the hypothesis in the sense required by the proof --- the theorem
is vacuous; under a literal reading, \eqref{eq:cond1} holds trivially for every
positive $\mathcal{A}/\mathcal{B}$, but the proof does not support the conclusion,
since already the constant $F_{r}$ alone satisfies $F_{r}\ge 9/5$ for $r\ge 6$ (for
instance $F_{6}=17017/9216>9/5$), so no choice of the tail primes can produce the
required contradiction. In either reading, no bound for $q_{r}$ with $r\ge 6$ is
obtained. We therefore suggest that the title, the abstract and the statement of
\cite[Theorem 1.2]{M10} be amended to refer to the $r$-th smallest prime divisor for
$2\le r\le 5$. Nothing else in \cite{M10} is affected: the alternative proof of
\cite[Theorem 1.1]{M10} and \cite[Corollary 1.1]{M10} use only $r\in\{2,3,4\}$.

Two further remarks on \cite{M10}. First, the range $2\le r\le 5$ does contain one
new bound that is not stated in \cite{M10}: for $r=5$ the threshold is
$1/(X_{5}-1)=5005/179$, so $(\mathcal{A},\mathcal{B})=(113,4)$ is admissible and the
theorem gives, in the sharper form $q_{r}<p_{L}$ established inside the proof,
\[
  q_{5}<p_{\lceil 113\,\omega(n)/4\rceil},
\]
in particular $q_{5}<p_{283}=1847$ when $\omega(n)=10$, the current lower bound for
$\omega(n)$ \cite{Thackeray}. Second, the reference ``From (3) and (4)'' in the last
display of the proof should be ``From (1) and (4)'', and the inequality
$\mathcal{A}>\mathcal{B}>1$ invoked in \cite[Remark 1.1]{M10} follows from
\eqref{eq:cond1} only in the range where $1/(X_{r}-1)>1$, which is again $2\le r\le5$.

\section{A misprint in Case 12 of \cite[Theorem 1.2]{CMM10}}

In Case 12 of the proof of \cite[Theorem 1.2]{CMM10} the chain under consideration is
$N=5^{2a_{1}}7^{2a_{2}}11^{2a_{3}}13^{2a_{4}}23^{2a_{5}}p_{6}^{2a_{6}}$, but the last
displayed factor in both abundancy products is $381/361=\Iab(19^{2})$ instead of
$\Iab(23^{2})=553/529$. With the correct factor the two inequalities still hold:
\[
  \Iab(5^{4}7^{2}11^{2}13^{2}23^{2})=\frac{1111642101}{614631875}>\frac95,
  \qquad
  \Iab(5^{2}7^{2}11^{2}13^{2}23^{2}31^{2})=\frac{15547332483}{8383578775}>\frac95 ,
\]
so the conclusion of Case 12, and hence \cite[Theorem 1.2]{CMM10}, is unaffected.

\section{Verification}\label{sec:verif}

Every numerical assertion above (the counterexamples to Lemma~\ref{lem:orig}; the
attainment of the minimum in Lemma~\ref{lem:min} for $1\le a\le 7$;
Proposition~\ref{prop:tight} checked by exhaustive minimisation over all admissible
configurations for $1\le a\le 6$ and $2a\le s\le 14$; the witnesses showing that
\eqref{eq:7} does not follow from the displayed inequality; the values of $X_{r}$ and
$F_{r}$ and the identity $F_{r}X_{r}=9/5$ for $2\le r\le 12$; and the two abundancy
products of Case 12) is verified by a deterministic script using exact integer and
rational arithmetic only, with an accompanying test suite. The code is available at
the repository cited in the author footnote.

\begin{thebibliography}{99}

\bibitem{CMM10} T.~Chatterjee, S.~Mandal and S.~Mandal,
\emph{A note on necessary conditions for a friend of $10$},
arXiv:2404.00624v5 (2025).

\bibitem{CMM20} T.~Chatterjee, S.~Mandal and S.~Mandal,
\emph{On characterizing potential friends of $20$},
Ann.\ West Univ.\ Timi\c{s}oara --- Math.\ Comput.\ Sci.\ \textbf{61} (2025), no.~1,
205--229; arXiv:2409.04451v4.

\bibitem{M10} S.~Mandal,
\emph{Prime divisors of $10$'s friends: a generalization of prior bounds},
An.\ Univ.\ Oradea Fasc.\ Mat.\ \textbf{33} (2026), no.~1, 5--12; arXiv:2412.02701v4.

\bibitem{M20div} S.~Mandal,
\emph{Exploring the relationships between the divisors of friends of $10$},
News Bull.\ Calcutta Math.\ Soc.\ \textbf{48} (2025), no.~1--3, 21--32;
arXiv:2504.08295v1.

\bibitem{MM} S.~Mandal and S.~Mandal,
\emph{Upper bounds for the prime divisors of friends of $10$},
Resonance \textbf{30} (2025), 263--275; arXiv:2404.05771v2.

\bibitem{Nielsen} P.~P.~Nielsen,
\emph{Odd perfect numbers, Diophantine equations, and upper bounds},
Math.\ Comp.\ \textbf{84} (2015), 2549--2567.

\bibitem{Thackeray} H.~R.~Thackeray,
\emph{Each friend of $10$ has at least $10$ nonidentical prime factors},
Indag.\ Math.\ \textbf{35} (2024), no.~3, 595--607; arXiv:2310.15900.

\bibitem{Ward} J.~Ward,
\emph{Does ten have a friend?}, Int.\ J.\ Math.\ Comput.\ Sci.\ \textbf{3} (2008),
no.~3, 153--158; arXiv:0806.1001.

\end{thebibliography}

\end{document}
